有限步正交化重分配矩阵更新在不同奇异方向上的能量，改变参数步长承担的曲率代价。给定相同梯度、动量和Hessian谱，
只需交换低曲率方向，就能改变线优化后的最优深度。二维相图给出了这一变化随梯度比和曲率比移动的精确分界。
对实际训练步长，方向曲率解释了采用较少迭代的候选为何会失去较高的一阶收益。
极因子误差此时也能准确排序，减小步长后则有更多状态由较浅候选取得最大下降。

在等幅训练中，三步增至五步带来的收益延续到较长训练。目标反馈的深度选择集中于NS5，
固定NS5在三组语言模型对照中取得较低的平均测试损失。
在实测步长下，低谱误差候选反复占优，在线探测多次返回同一个深度。
对发生切换的两条轨迹，采用NS3的时机解释了反馈与固定深度之间的大部分测试损失差距。
